% !TeX root = prueba01_version01.tex
% !TeX program = pdflatex
% Propuesta parcial de la primera evaluación.
\documentclass[10pt,a5paper]{article}
\newif\ifmostrarSoluciones
\mostrarSolucionestrue
\usepackage{estilo_apuntes}
\dosPorCarillafalse % Cambie a \dosPorCarillatrue para imprimir dos páginas por carilla.
\encabezadoprueba{Prueba 1}{22 de octubre de 2026}

\begin{document}
\begin{center}
  {\color{Azul}\large\bfseries Prueba 1}
\end{center}

\medskip

\medskip
\noindent\textbf{Indicaciones}
\begin{itemize}
  \item La duración de la prueba es de 1 hora y 30 minutos.
  \item Escriba el desarrollo de cada ejercicio de manera clara y ordenada.
  Justifique cada paso e indique las propiedades, restricciones o resultados que
  utilice. Una respuesta sin justificación recibirá únicamente crédito parcial.
  \item No se permite el uso de ningún dispositivo electrónico durante la
  evaluación.
  \item Cualquier intento de copia penalizará la evaluación con una calificación
  total de cero.
\end{itemize}

\medskip
\noindent\textbf{Ejercicios}
\begin{enumerate}
  \item Sean $a,b\in\mathbb{R}$ y sea $F\colon\mathbb{R}\to\mathbb{R}$ la función definida por
  \[
    F(x)=
    \begin{cases}
      \dfrac{x^2-1}{x-1}, & x<1,\\[2mm]
      5, & x=1,\\
      ax+b, & x>1.
    \end{cases}
  \]
  Determine $a$ y $b$ para que exista
  \[
    \lim_{x\to1}F(x)
  \]
  y, además, se cumpla $F(2)=5$. Calcule el límite.

  \item Sea $B$ una función definida para $x$ cerca de $0$, excepto posiblemente en $0$, que satisface
  \[
    |B(x)|\leq4.
  \]
  Calcule mediante el teorema del Sanduche
  \[
    \lim_{x\to0}
    \frac{\sqrt{1+x^2}-1}{x}\,B(x).
  \]
  \textbf{Hint:} Note que $\sqrt{1+x^2}\geq 1$.
  \item Considere las funciones
  \[
    f\colon[1,\infty)\to[0,\infty),
    \qquad
    f(x)=\sqrt{x-1},
  \]
  y
  \[
    g\colon\mathbb{R}\to[1,\infty),
    \qquad
    g(x)=x^2+1.
  \]
  \begin{enumerate}
    \item Determine si cada función es inyectiva, sobreyectiva o biyectiva.
    \item Determine la función $g\circ f$.
  \end{enumerate}
\end{enumerate}

\ifmostrarSoluciones
\begin{solucion}[de la propuesta]
\begin{enumerate}
  \item Para $x<1$ y $x\neq1$,
  \[
    \frac{x^2-1}{x-1}=x+1.
  \]
  Por ello,
  \[
    \lim_{x\to1^-}F(x)=2.
  \]
  Por la derecha,
  \[
    \lim_{x\to1^+}F(x)=a+b.
  \]
  Como ambos límites laterales son finitos, el límite bilateral existe si y
  solo si son iguales, es decir, si y solo si
  \[
    a+b=2.
  \]
  Además, la condición $F(2)=5$ implica
  \[
    2a+b=5.
  \]
  Al resolver el sistema se obtiene
  \[
    a=3,
    \qquad
    b=-1.
  \]
  En consecuencia,
  \[
    \lim_{x\to1}F(x)=2.
  \]

  \item Para $x\neq0$, racionalizamos el factor algebraico:
  \[
    \frac{\sqrt{1+x^2}-1}{x}
    =\frac{x}{\sqrt{1+x^2}+1}.
  \]
  Entonces
  \begin{align*}
    \left|
      \frac{\sqrt{1+x^2}-1}{x}B(x)
    \right|
    &=\frac{|x|}{\sqrt{1+x^2}+1}|B(x)|\\
    &\leq\frac{4|x|}{\sqrt{1+x^2}+1}\\
    &\leq2|x|.
  \end{align*}
  Por tanto,
  \[
    -2|x|
    \leq
    \frac{\sqrt{1+x^2}-1}{x}B(x)
    \leq
    2|x|.
  \]
  Como ambas cotas tienden a $0$, el teorema del Sanduche permite concluir que
  \[
    \lim_{x\to0}
    \frac{\sqrt{1+x^2}-1}{x}B(x)=0.
  \]

  \item La función $f$ es estrictamente creciente en $[1,\infty)$, por lo que es
  inyectiva. Además, para cada $y\in[0,\infty)$, el punto
  \[
    x=y^2+1
  \]
  pertenece a $[1,\infty)$ y satisface $f(x)=y$. Por tanto, $f$ es sobreyectiva
  y, en consecuencia, biyectiva.

  La función $g$ es sobreyectiva sobre $[1,\infty)$: dado $y\geq1$, basta tomar
  \[
    x=\sqrt{y-1}.
  \]
  Sin embargo, no es inyectiva porque, por ejemplo,
  \[
    g(1)=g(-1)=2.
  \]
  Por ello, $g$ no es biyectiva.

  Como $f(x)\in[0,\infty)\subseteq\mathbb{R}$, la composición $g\circ f$ está
  definida para todo $x\in[1,\infty)$ y
  \[
    (g\circ f)(x)
    =\left(\sqrt{x-1}\right)^2+1
    =x.
  \]
  Por tanto, la función solicitada queda definida por
  \[
    g\circ f\colon[1,\infty)\to[1,\infty),
    \qquad
    (g\circ f)(x)=x.
  \]
  Es la función identidad en $[1,\infty)$.
\end{enumerate}
\end{solucion}
\fi

\end{document}
